\begin{titlepage}
\centering
\vspace*{2in}
{\Huge A First Course in\\[0.3em] Monte Carlo Methods\par}
\vspace{1.5em}
{\large D.\ Sanz-Alonso and O.\ Al-Ghattas\par}
\vspace{0.5em}
{\normalsize University of Chicago\par}
\vfill
\end{titlepage}

\frontmatter

\listoffigures
\listofalgorithms
\addcontentsline{toc}{chapter}{List of Algorithms}

\chapter*{Preface}
\addcontentsline{toc}{chapter}{Preface}
\markboth{Preface}{Preface}

\textbf{Overview}\quad This textbook provides a concise mathematical introduction to Monte Carlo
methods, a rich family of algorithms with far-reaching applications in science and engineering.
Monte Carlo methods are an exciting topic for computational mathematicians and statisticians:
the design and analysis of modern algorithms are rooted in a broad mathematical toolbox that
includes ergodic theory of Markov chains, Hamiltonian dynamical systems, transport maps,
stochastic differential equations, information theory, optimization, Riemannian geometry, and
gradient flows, among many others. This concise textbook celebrates the breadth of mathematical
ideas that have led to tangible advancements in Monte Carlo methods and their applications.
To accommodate a broad audience, the level of mathematical rigor varies across chapters, with
the most technically demanding subjects receiving only an intuitive treatment.

The aim is not to be comprehensive or encyclopedic, but rather to illustrate some key principles
in the design and analysis of Monte Carlo methods through a carefully-crafted choice of topics
that emphasizes timeless over timely ideas. Algorithms are presented in a way that is conducive
to conceptual understanding and mathematical analysis---clarity and intuition are favored over
state-of-the-art implementations that are harder to comprehend or rely on \textit{ad-hoc}
heuristics. To help readers navigate the expansive landscape of Monte Carlo methods, each
algorithm is accompanied by a discussion of its strengths and limitations, as well as the types
of problems for which it is most suitable. The presentation is self-contained, and therefore
adequate for self-guided learning and as a teaching resource. Each chapter includes a section
with bibliographic remarks to guide computational scientists and graduate students interested
in research on Monte Carlo methods and their applications.

\vspace{1em}
\textbf{Acknowledgements}\quad This concise textbook developed as a result of several courses
taught by Daniel Sanz-Alonso at the University of Chicago since 2019; Omar Al-Ghattas served
as teaching assistant for two of these courses and contributed significantly to improve the
presentation. The main target audience are advanced undergraduate and graduate students in
computational mathematics, applied mathematics, and statistics.

\hfill D.\ Sanz-Alonso and O.\ Al-Ghattas

\hfill University of Chicago

\tableofcontents

\mainmatter
